\section{The calculus and the two comparison questions}\label{sec:calculus}

The ground type is $o$.  Simple types are formed by $\sigma\to\tau$,
with arrows associated to the right, and their levels are
\[
 \lv(o)=0,\qquad
 \lv(\sigma\to\tau)=\max\{\lv(\sigma)+1,\lv(\tau)\}.
\]
Thus $o\to o$ has level one, whereas $(o\to o)\to o$ has level two.
A type level measures the ability to take functions as arguments; it does
not measure the number of arguments.  A curried type $o^k\to o$ has
level one for every positive finite $k$.

The original calculus has variables, application and abstraction, together
with zero, successor, case functionals and recursors.  We use the types
\[
 0:o,\quad S:o\to o,\quad D^\tau:o\to\tau\to\tau\to\tau,
 \quad R^\tau:o\to(o\to\tau\to\tau)\to\tau\to\tau.
\]
In particular the order of the arguments to a recursor is fixed by
\[
 R^\tau0AB\longrightarrow B,\qquad
 R^\tau(Sn)AB\longrightarrow A\,n\,(R^\tau nAB).
\]
The case rules are $D^\tau0AB\to A$ and $D^\tau(Sn)AB\to B$.
Together with beta reduction these rules act inside arbitrary well-typed
contexts.  Terms in this paper are literal terms: an input is not first
beta-normalized, uncurried, or replaced by an extensionally equal program.

A termination ordering is a relation that decreases on these reductions
and is stable under contexts and substitutions.  Recursive path ordering
provides the corresponding first-order example.  For ranked trees,
$f(s_1,\ldots,s_k)>t$ follows either from a child dominating $t$, or from
domination of all children of $t$ together with a decrease of root precedence
or a strict multiset decrease between equal roots.  A strict multiset
decrease replaces a nonempty collection by elements each smaller than a
removed element.  Extending this mechanism to typed application and
abstraction introduces both type and binding information.

The separate ordinal-vector comparison studied below compares numerical
assignments under \emph{all} permitted substitutions.  It is not the same
definition as a recursive syntactic ordering.  Nor is the height of a
termination ordering the largest ordinal written in an assignment.
The Bachmann--Howard calibration concerns the descending height of the
syntactic relation.  Its named ordinal notation systems and their comparison
are supplied by the published ordinal analyses \cite{pohlers,simmons,simmonshelpful};
the exact new ordering and its proof are given in \cite{wwordering}.

\section{Ordinal expressions and the fixed numerical control}\label{sec:ordinals}

Every ordinal below $\varepsilon_0$ has a finite recursive Cantor normal
form.  Here $\varepsilon_0$ is the least fixed point of
$\alpha\mapsto\omega^\alpha$.  We write $\oplus,\otimes$ for natural
sum and product: collect the Cantor monomials for addition, and multiply
monomials by naturally adding their exponents for multiplication.  These
operations are commutative.  Ordinary ordinal addition is retained when
a Cantor normal form or a finite tail is displayed.  For example
$\omega\oplus1=\omega+1$ but $1+\omega=\omega$ for ordinary addition.

\begin{definition}[Additive tree norm]\label{def:norm}
For a Cantor normal form with decreasing exponents and positive coefficients,
put
\[
 \norm(0)=0,\qquad
 \norm\left(\sum_{i<r}\omega^{\alpha_i}c_i\right)
       =\sum_{i<r}c_i(1+\norm(\alpha_i)).
\]
Thus $\norm(n)=n$, $\norm(\omega)=2$ and
$\norm(\omega^{\omega+1})=4$.  Coefficients are charged with their full
multiplicity, not their binary length.
\end{definition}

\begin{lemma}[Finite norm balls]\label{lem:finite-balls}
Each norm ball is finite and can be effectively enumerated.  Cantor-normal-form
comparison, natural sum and natural product are effective.  Moreover
$\norm(\alpha\oplus\beta)=\norm(\alpha)+\norm(\beta)$.
\end{lemma}
\begin{proof}
An ordinal of norm at most $b$ is represented by a finite rooted expression
tree with at most $b$ charged monomial nodes, including repeated monomials.
There are finitely many such finite trees.  Enumerate them, interpret them
recursively, and retain their canonical Cantor forms.  Alternatively enumerate
by norm recursively: every exponent in a nonzero term of norm at most $b$
has norm at most $b-1$, so induction on $b$ gives a finite list.
Comparison is lexicographic comparison of exponent/coefficient lists, using
the same recursive algorithm on exponents.  The definitions of natural sum
and product then give finite algorithms.  Collecting repeated monomials does
not change the sum of their charged norms, proving additivity.
\end{proof}

Fix throughout
\begin{equation}\label{eq:phi-definition}
 F_0(t)=2^t,\qquad F_{j+1}(t)=F_j^{\,t+1}(t),\qquad
 \Phi(t)=F_5(t+100),
\end{equation}
where the exponent denotes iteration.  Replacing the exponential base
function by doubling would change every subsequent definition.

\begin{lemma}[Elementary control properties]\label{lem:phi-growth}
The fixed $\Phi$ is computable, strictly increasing, and greater than the
identity.  On its whole domain,
\[
 \Phi(t+1)\ge 2^{\,2^{\Phi(t)}}.
\]
In particular $\Phi(t+1)>4\Phi(t)+4$.
\end{lemma}
\begin{proof}
Induction on $j$ proves that each $F_j$ is computable, increasing, and
strictly above the identity.  For $n\ge1$ the same induction gives
$F_{j+1}(n)\ge F_j(n)$ and $F_1(n)\ge2^{2^n}$.  The $n+1$ evaluations
inside $F_5(n+1)$ start from a strictly larger seed than those in $F_5(n)$.
Since all values are integers,
\[
 F_5(n+1)=F_4(F_4^{n+1}(n+1))\ge F_4(F_5(n)+1)
             \ge2^{2^{F_5(n)+1}}.
\]
Take $n=t+100$.  The stated weaker inequality follows.  Its final consequence
holds because $\Phi(t)\ge4$ and $2^{2^a}>4a+4$ for $a\ge4$.
\end{proof}

\begin{definition}[Norm-controlled collapse]\label{def:collapse}
For $\alpha<\varepsilon_0$ define, by ordinal recursion,
\[
 \psi(\alpha)=\max\bigl(\{0\}\cup
 \{\psi(\beta)+1:\beta<\alpha,\ %
                 \norm(\beta)\le\Phi(\norm(\alpha))\}\bigr).
\]
The set is finite by Lemma~\ref{lem:finite-balls}.
\end{definition}

\begin{lemma}[Controlled strictness and finite values]\label{lem:controlled-strictness}
If $\beta<\alpha$ and
$\norm(\beta)\le\Phi(\norm(\alpha))$, then
$\psi(\beta)+1\le\psi(\alpha)$.  Also $\psi(n)=n$ for finite $n$.
\end{lemma}
\begin{proof}
The first assertion is one term of the defining maximum.  For the second,
induct on $n$.  Every predecessor of the finite ordinal $n$ is finite and
has norm below $n\le\Phi(n)$.  The induction hypothesis makes the maximum
equal to $n$.
\end{proof}

The further inequalities $\norm(\alpha)\le\psi(\alpha)$ and
$k\le\psi(\alpha+k)$, as well as the natural-product norm bounds used
in source assignments, are Proposition~2.3 and Proposition~2.2 of
Wilken--Weiermann \cite{ww}.  Those are published background theorems,
not assumptions about a new comparison algorithm.  Controlled strictness
does not assert that $\psi$ is globally monotone in ordinal order.

\section{A complete closed evaluator}\label{sec:evaluator}

\begin{definition}[Maximum predecessor]\label{def:maximum}
For $\alpha>0$ and $b\in\N$ let
\[
 M(\alpha,b)=\max\{\beta<\alpha:\norm(\beta)\le b\},\qquad
 L(e,b)=M(\omega^e,b).
\]
The maximum exists because the set is finite and contains zero.
\end{definition}

\begin{theorem}[Effective maximum predecessor]\label{thm:maximum}
The following simultaneous recursion computes $M$ and $L$ exactly.
Write $\alpha=\omega^e c+r$ with $c>0$ and $r<\omega^e$, and put
$w=1+\norm(e)$.
\begin{enumerate}
\item $M(\alpha,0)=0$ and $L(0,b)=L(e,0)=0$.
\item If $r>0$ and $cw\le b$, retain the leading block:
 $M(\alpha,b)=\omega^e c+M(r,b-cw)$.
\item Otherwise set $d=\min(c-1,\lfloor b/w\rfloor)$ when $r=0$,
 and $d=\min(c,\lfloor b/w\rfloor)$ when $r>0$.  Then
 $M(\alpha,b)=\omega^e d+L(e,b-dw)$, omitting a zero block.
\item For $e,b>0$, let $a=M(e,b-1)$, $v=1+\norm(a)$ and
 $(q,s)=\operatorname{divmod}(b,v)$.  Then
 $L(e,b)=\omega^a q+L(a,s)$.
\end{enumerate}
\end{theorem}
\begin{proof}
If a nonzero tail can be decreased while retaining the entire leading block,
every predecessor with a smaller leading block is inferior.  The remaining
budget and tail bound give clause~2.  Otherwise the displayed $d$ is the
largest affordable coefficient still strictly below the original bound.
After choosing that coefficient the tail must lie below $\omega^e$, which
is exactly the $L$ query in clause~3.

For clause~4, any nonzero candidate below $\omega^e$ of norm at most $b$
has leading exponent $a'<e$ with $\norm(a')\le b-1$.  The largest possible
exponent is therefore $a=M(e,b-1)$.  Lexicographic Cantor comparison first
maximizes that exponent, then its coefficient $q$, then the remaining tail.
The last tail is bounded by $\omega^a$ and has budget $s$, proving the formula.

For termination use the lexicographic pair (budget, phase), putting $M$
above $L$ at equal budget.  Clause~2 lowers the budget since $cw>0$.
Clause~3 may keep it but lowers the phase.  In clause~4, the $M$ call uses
$b-1$, and $s<b$ because $v\le b$ and hence $q\ge1$.  Base cases terminate.
This proves a finite algorithm and the asserted exactness, without an oracle
for the maximum.
\end{proof}

\begin{lemma}[Saturation at a limit]\label{lem:saturation}
For every nonzero limit $\lambda$ and $b\in\N$,
$\norm(M(\lambda,b))=b$.  The maximum strictly increases with $b$.
\end{lemma}
\begin{proof}
If $\eta=M(\lambda,b)$ left unused budget $d>0$, then $\eta+d<\lambda$
would have norm $b$ and be larger.  This contradicts maximality.  For $b<c$,
the ordinal $M(\lambda,b)+1$ is below the limit and has norm $b+1\le c$.
It witnesses strict increase.
\end{proof}

\begin{theorem}[Single-predecessor recurrence]\label{thm:single-predecessor}
For every nonzero limit $\lambda<\varepsilon_0$ and $n\in\N$,
\[
 \psi(\lambda+n)=
 \psi\bigl(M(\lambda,\Phi(\norm(\lambda)+n))\bigr)+1.
\]
\end{theorem}
\begin{proof}
Induct on $n$.  Put $b=\Phi(\norm(\lambda)+n)$ and
$\eta=M(\lambda,b)$.  This is an admitted predecessor and has norm $b$,
so its contribution gives the lower inequality.  Any other admitted
$\beta<\lambda$ satisfies $\beta\le\eta$.  For $\beta<\eta$ the norm
premise $\norm(\beta)\le b=\norm(\eta)\le\Phi(\norm(\eta))$ and
Lemma~\ref{lem:controlled-strictness} bound its contribution by $\psi(\eta)$.

The remaining predecessors have the form $\lambda+j$, $j<n$.  Their budgets
$b_j=\Phi(\norm(\lambda)+j)$ are strictly smaller than $b$.  By saturation
$\eta_j=M(\lambda,b_j)<\eta$, with $\norm(\eta_j)=b_j<b$.
Controlled strictness and the induction hypothesis give
$\psi(\lambda+j)=\psi(\eta_j)+1\le\psi(\eta)$.  Every term in the
original maximum is therefore at most $\psi(\eta)+1$.
\end{proof}

\begin{corollary}[Closed evaluation]\label{cor:closed-evaluation}
Every supplied closed collapse value is computable.  At a nonfinite state
$\alpha=\lambda+n$, replace $\alpha$ by
$M(\lambda,\Phi(\norm(\alpha)))$ and increment a counter.  On reaching
a finite state return that state plus the counter.
\end{corollary}
\begin{proof}
Each replacement strictly decreases the ordinal.  Well-foundedness below
$\varepsilon_0$ proves termination.  The recurrence proves preservation
of the desired value, and Lemma~\ref{lem:controlled-strictness} evaluates the
last finite state.  This is pointwise computability, not a uniform practical
complexity bound and not a decision procedure for an infinite family of inputs.
\end{proof}

\begin{theorem}[Complete terminal decomposition]\label{thm:complete-count}
Put $\nu=\norm(\lambda)$ and let $t\ge\nu$ be the total initial norm.
Starting at $\lambda+(t-\nu)$, let $\tau_\lambda(t)$ be the number of
nonfinite replacements and $Z_\lambda(t)$ the first finite state.  Then
\[
 J_\lambda(t):=\psi(\lambda+t-\nu)
       =Z_\lambda(t)+\tau_\lambda(t),\qquad
 Z_\lambda(t)=\Phi^{\tau_\lambda(t)}(t).
\]
The count is positive.  In particular $Z_\lambda(t)<J_\lambda(t)<2Z_\lambda(t)$.
\end{theorem}
\begin{proof}
Saturation makes every new norm exactly $\Phi$ of the old norm, including
the last replacement into a finite state.  Each replacement contributes
one to the recurrence.  At the finite endpoint the value equals its norm.
There is no subsequent count of its finite predecessors.
Since each application of $\Phi$ increases its argument by at least one,
$Z_\lambda(t)\ge t+\tau_\lambda(t)$.  A nonzero limit has $t\ge2$;
therefore $0<\tau_\lambda(t)\le Z_\lambda(t)-2$ and the strict bounds follow.
\end{proof}

\begin{corollary}[Finite coefficient identity]\label{cor:finite-head}
For finite $k,n$,
\[
 \psi(\omega k+n)=\Phi^k(2k+n)+k.
\]
\end{corollary}
\begin{proof}
The case $k=0$ is finite evaluation.  For $k>0$, with
$b=\Phi(2k+n)\ge2k$, the maximum predecessor is
$\omega(k-1)+b-2(k-1)$.  Apply the recurrence and induction on $k$.
\end{proof}

\begin{lemma}[Strict separation in one finite tail]\label{lem:tail-separation}
For $h_\lambda(n)=\psi(\lambda+n)$ with a fixed nonzero limit $\lambda$,
\[
 h_\lambda(n+1)>2^{h_\lambda(n)}.
\]
In particular every such function is injective.
\end{lemma}
\begin{proof}
Compare descents at total norms $t$ and $t+1$, using $t\ge\norm(\lambda)$.
Maximum predecessors are monotone in their ordinal bound and norm budget;
dropping a finite tail also preserves ordinal order.  Before the smaller
descent reaches a finite state, induction on the simultaneous step number
shows that the larger descent has not reached a finite state either.  Its
count $L'$ is at least the smaller count $L\ge1$.
Writing $Z=\Phi^L(t)$, monotonicity and Lemma~\ref{lem:phi-growth} give
\[
 \Phi^{L'}(t+1)\ge\Phi^L(t+1)\ge 2^{2^Z}.
\]
For the last inequality, apply the control inequality to the last iteration;
its two input integers differ by at least one.  By the preceding theorem
$h_\lambda(n)=Z+L<2Z\le2^Z$.  The larger value is at least its endpoint,
which proves the strict inequality.
\end{proof}

The maximum-predecessor framework is classical \cite{bcw}.  The explicit
recursion above fixes its exact specialization to this norm and collapse.
The recurrence also holds for any computable strictly increasing control
majorizing the identity.  The very strong tail inequality uses the particular
control \eqref{eq:phi-definition}.

\section{The canonical source assignment}\label{sec:assignment}

For $A:\tau$ the canonical vector has indices $0,\ldots,\lv(\tau)$.
Missing indices are padded by zero when vectors are compared.  The prime
clauses are
\[
 \can0=(1),\quad \can S=(1,1),\quad
 \can{D^\tau}=(1,\ldots,1),\quad
 \can{R^\tau}=(2,1,\ldots,1,\omega),\quad \can X=\mathbf x.
\]
For vector levels $m>n$, compute $h=f\square g$ downward by
\[
 h_i=f_i\ (i>n),\qquad
 h_i=2^{h_{i+1}}\otimes(f_i\oplus g_i)\ (1\le i\le n),\qquad
 h_0=\psi(\omega\otimes h_1\oplus f_0\oplus g_0\oplus n).
\]
Application retains the coordinates through the result type's level.
In the literal scalar tree we associate the tail as $f_0+(g_0+n)$ for
$n>0$ and as $f_0+g_0$ for $n=0$.  This is relevant to syntactic partials.

For abstraction over $X:\sigma$, put $n=\lv(\sigma)+1$ and let $h$
have level $m$.  Each auxiliary partial $\partial_i^Xh_i$ is a vector
of level $n$.  The canonical zero-history rule is
\[
 \can{\lambda X.G}=\delta^X\can G+\can G\{X:=\mathbf1\},
\]
where
\[
 (\delta^Xh)_0=2^n\left(\sum_{i=0}^m\sz^X(h_i)\right)(\partial_0^Xh_0)_0,
\quad
 (\delta^Xh)_j=\sum_{i=0}^m(\partial_i^Xh_i)_j\ (0<j\le n),
\]
and $(\delta^Xh)_j=h_j$ for $j>n$.
The size is one at a whole $X$-free expression or variable leaf; an active
sum or collapse adds the child sizes and one; $2^f\otimes g$ has size
$2\sz^X(f)+\sz^X(g)+1$.

At a whole $X$-free input $u$, $(\partial_i^Xu)_j=u+1$ if $i=j$ and
is one otherwise.  At $x_i$ every output is one.  Active sums add the
partial vectors and $\mathbf1$.  For $i>0$,
\[
 \partial_i^X(2^f\otimes g)=2\partial_{i+1}^Xf+\partial_i^Xg+\mathbf1.
\]
At an active scalar $\psi(\omega\otimes f+g)$, the zeroth output is
$\psi(\omega\otimes f\{X:=\mathbf1\}+(\partial_0^Xg)_0)$, and output
$j>0$ is $(\partial_1^Xf)_j+(\partial_0^Xg)_j$.
The whole-free rule applies before any other rule, on the literal tree.
These are the published clauses of \cite[Definitions 2.10, 2.18 and 3.1]{ww}.
Their domain and norm-preservation theorems are imported from that paper.

The closed graded classes $\mathcal B_i$ are defined simultaneously as follows:
one belongs to every class; $\omega$ belongs to classes of positive index;
each class is closed under natural sum; $2^f\otimes g\in\mathcal B_i$
when $i\ge1$, $f\in\mathcal B_{i+1}$ and $g\in\mathcal B_i$;
$\psi(\omega\otimes f+g)\in\mathcal B_0$ when
$f\in\mathcal B_1$, $g\in\mathcal B_0$ and $\norm(f)\le F_2(g)$;
and a member of $\mathcal B_0$ may also be used in any positive class.
These are classes of closed expression trees; zero is not a leaf of the
grammar, although zero padding is used outside a vector's actual length.
Admissible substitutions replace each entire variable vector by a closed
vector with $x_i\in\mathcal B_i$ and $\norm(x_i)\le x_0$ for every
positive index.  A ground variable therefore ranges over all positive
integers, independently of other names: every such integer is a finite sum
of ones and there are no positive coordinates imposing a bound.
Weak comparison is coordinatewise universal comparison.  Strict comparison
is strict at coordinate zero and weak at every positive coordinate; it is
not lexicographic.

\begin{theorem}[Actual-source forward translation]\label{thm:forward-translation}
For a literal original source with only ground free names, all positive
canonical coordinates are closed and effectively evaluable.  Coordinate
zero is effectively a finite expression in positive-natural inputs, constants,
addition and fixed functions $h_\lambda(n)=\psi(\lambda+n)$, where each
direct index is $\lambda=\omega\otimes c>0$ for a closed $c$.
The translation preserves every shared input valuation, including arbitrary
internal types, binders and recursors.
\end{theorem}
\begin{proof}
Use simultaneous induction on source construction and the syntactic partial
recursions.  The invariant for a generated positive coordinate is that its
parse uses positive-index variables, positive constants, natural addition
and the graded exponential/product constructor, but no ground-index variable.
The prime clauses have this property.  The box computes positive fields
using positive fields only and preserves it.

For a positive input to a partial, the whole-free rule retains only that
positive input at its diagonal output and constants elsewhere.  Thus it
does not introduce a ground variable.  At an active variable the result is
constant.  Sum and exponential clauses preserve the invariant recursively.
For a scalar input, consider a positive output of its partial: the whole-free
rule returns one, a tracked scalar variable returns one, and sums recurse.
The collapse clause uses a positive-head partial and the positive output of
the scalar-tail partial.  These are exactly the two induction hypotheses.
Hence even partials of scalar expressions cannot introduce ground-index
variables into their positive outputs.  The abstraction sums such outputs
and substitutes ones for the bound vector in its retained part, so it
preserves the invariant and removes all of that bound name's components.

Positive fields of the final vector can therefore mention only positive
components of free term variables.  There are none when every free variable
is ground.  Closed evaluation follows from the ordinal algorithms already
proved.  In particular the head of every final scalar collapse is closed.

At coordinate zero the prime and box clauses produce only constants,
ground inputs, addition and collapse nodes.  The scalar partial clause has
the same form, and its size multiplier is a fixed syntactic integer, not a
function of a valuation.  Expanding it as a finite multiplicity preserves
the grammar.  Replace each closed-head collapse by $h_{\omega\otimes c}$.
All replacements are compositional and use the original free names, proving
the exact, valuation-preserving translation.

This argument concerns generated assignments.  The ambient graded classes
also allow promotion of certain whole-free scalar expressions to positive
levels; the stronger invariant just proved is not asserted for all arbitrary
members of those ambient classes.
\end{proof}

\begin{corollary}[Two natural positive classes]\label{cor:source-classes}
In a recursor-free source every positive coordinate is finite.  If every
internal type and subtype has level at most two, every direct scalar index
is a polynomial limit below $\omega^\omega$.
\end{corollary}
\begin{proof}
Only the recursor prime introduces $\omega$ into a generated positive field.
All other positive constructors, including the positive partial cases in
the preceding proof, preserve finiteness after the final closures.  This
proves the first assertion for arbitrary internal type levels.
For the second, coordinate two cannot contain an exponential constructor
requiring coordinate three.  Its generated closed values are linear Cantor
polynomials $a\omega+b$.  Coordinate one is built by sums and natural products
using $2^{a\omega+b}=\omega^a2^b$, hence lies in $\N[\omega]$.
Multiplication naturally by $\omega$ leaves a polynomial limit.  A bound
only on the output type would not bound internal coordinates in this way.
\end{proof}

The converse image assertion is not part of the translation: not every
abstract expression need have an original source realizing it.  Constructing
such sources, retaining their finite box and input correlations, is one of
the separate tasks addressed in \cite{wwexpressivity}.

\subsection{A worked example}
Let $I=\lambda z:o.z$ and $K=\lambda z:o.0$.  The partial and size rules
give $\can I=(3,1)$ and $\can K=(5,1)$: in the first case the scalar
partial is one, in the second it is $1+1$, and the size multiplier is two.
Thus, for a raw ground input $x_0>0$,
\[
 \can{Sx}_0=\Phi(x_0+3)+1,\quad
 \can{Ix}_0=\Phi(x_0+5)+1,\quad
 \can{Kx}_0=\Phi(x_0+7)+1.
\]
These are assigned values, not program denotations.  Although $Ix\to x$,
the canonical values differ.  Canonical assignments fix the zero-history
choice at every lambda; the normalization theorem in \cite{ww} can use
history-dependent assignments.  One cannot replace the input syntax by its
normal form when deciding its canonical comparison.
